\subsection{CASE OF A COMMUTATIVE FIELD}

All the above applies when the ring A is a commutative field; but there are simplifications and additional results.

Thus Proposition 12 of § 7, no. 9 can be formulated in this case as follows:

PROPOSITION 15. Let E be a vector space over a commutative field; for p vectors \(x_{i} \in E\) ( \(1 \leqslant i \leqslant p\) ) to be linearly independent, it is necessary and sufficient that \(x_{1} \wedge x_{2} \wedge \cdots \wedge x_{p} \neq 0\) .

COROLLARY. Let X be a matrix of type \((m, n)\) over a commutative field. The rank of X is equal to the greatest integer p such that there exists at least one minor of X of order p which is \(\neq0\) .

The rank of \(X\) is the maximum number of linearly independent columns of \(X\) (II, § 10, no. 12, Definition 7). The corollary then follows from Proposition 15 and formula (17) of no. 5.

Consider now the case of a system of m linear equations in n unknowns over a commutative field K:

\[
\sum_ {j = 1} ^ {n} \lambda_ {i j} \xi_ {j} = \eta_ {i} \quad (1 \leqslant i \leqslant m).\tag{41}
\]

PROPOSITION 16. Let \(L = (\lambda_{ij})\) be the matrix (of type \((m, n)\) ) of system (41). Let \(M\) be the matrix of type \((m, n + 1)\) obtained by adding to \(L\) the \((n + 1)\) -th column \((\eta_i)\) (II, § 10, no. 1). Let \(p\) be the rank of \(L\) (calculated by applying the Corollary to Proposition 15). Suppose that the minor \(\Delta\) of \(L\) , the determinant of the matrix by suppressing the rows and columns of index \(\geqslant p + 1\) in \(L\) , is \(\neq 0\) (which is always possible by means of a suitable permutation on the rows of \(L\) and a suitable permutation on the columns of \(L\) ). Then, for system (41) to have at least one solution it is necessary and sufficient that all the minors of order \(p + 1\) of \(M\) , the determinants of the submatrices of order \(p + 1\) of \(M\) whose columns have indices \(1, 2, \ldots, p\) and \(n + 1\) , be zero. If this is so, the solutions of system (41) are those of the system consisting of the first \(p\) equations; if they are written

\[
\sum_ {j = 1} ^ {p} \lambda_ {i j} \xi_ {j} = \eta_ {i} - \sum_ {k = p + 1} \lambda_ {i k} \xi_ {k} \quad (1 \leqslant i \leqslant p)\tag{42}
\]

all the solutions of this system are obtained by taking for the \(\xi_{k}\) of index \(k > p\) arbitrary values and applying Cramer's formulae (no. 7, formulae (37)) to calculate the \(\xi_{j}\) of index \(j \leqslant p\) .

We know (II, § 10, no. 12, Proposition 12) that for the system (41) to have at least one solution, it is necessary and sufficient that the matrices L and M have the same rank. With the rows and columns of L permuted to satisfy the condition of the statement, let \(a_{i}\) ( \(1 \leqslant i \leqslant p\) ) denote the first p columns of L and \(y = (\eta_{i})\) the \((n + 1)\) -th column of M; since all the columns of L are by hypothesis linear combinations of the \(a_{i}\) , to say that M has the same rank p as L means that y is a linear combination of the \(a_{i}\) , or also (Proposition 15) that \(a_{1} \wedge \cdots \wedge a_{p} \wedge y = 0\) . The possibility condition in the statement is the translation of the latter relation, taking account of formula (17) of no. 5. Moreover, since the first p rows of M are linearly independent, the rows of index \textgreater p are linear combinations of them and hence every solution of (42) is also a solution of (41). The last assertion is then an immediate consequence of Proposition 13 of no. 7.

